\section{Introduction and Motivation}
\label{sec:intro}

\subsection{Reduced-Order Modeling and Operator Learning for Time-Dependent PDEs}

The numerical approximation of parametrized, time-dependent partial differential equations (PDEs) by high-fidelity full-order models (FOMs)---such as finite element, finite volume, or spectral discretisations---often entails state spaces of dimension $N_h \sim 10^4$ to $10^8$. When simulations must be repeated across hundreds of input signals, boundary controls, or physical parameters, as in optimal control, uncertainty quantification, or real-time digital twins, the computational cost of querying the full-order solver becomes prohibitive. Classical projection-based reduced-order models (ROMs), built upon Proper Orthogonal Decomposition (POD) and Galerkin or Petrov--Galerkin projection~\cite{quarteroni2015reduced,hesthaven2016certified}, address this bottleneck by approximating the solution field $y(\mathbf{x},t)$ within a low-dimensional linear trial subspace $V_K = \operatorname{span}\{\varphi_1(\mathbf{x}),\dots,\varphi_K(\mathbf{x})\}$.

Despite their rigorous error theory and physical structure preservation on elliptic and parabolic problems, linear subspace ROMs face two well-known structural obstacles. First, when the PDE is non-affine or strongly nonlinear, assembling the projected reduced operator requires hyper-reduction techniques and remains intrusive with respect to the full-order solver code. Second, and more fundamentally, any method that approximates the solution in a fixed $K$-dimensional linear subspace of $L^2(\Omega)$ is bounded below in accuracy by the Kolmogorov $n$-width of the solution manifold~\cite{pinkus1985n,ohlberger2016reduced}. For diffusion-dominated problems or smooth vortex flows, the singular values of the snapshot matrix decay exponentially fast, so a dozen modes suffice to reach sub-percent error. By contrast, for advection-dominated problems, hyperbolic conservation laws, or reaction--diffusion equations exhibiting sharp travelling wave fronts (such as cardiac action potentials~\cite{aliev1996simple}), translates of a sharp front of width $\delta \ll 1$ are nearly orthogonal in $L^2(\Omega)$. Consequently, the Kolmogorov $n$-width decays slowly ($O(K^{-1/2})$ for discontinuous shocks, or with a very slow pre-asymptotic geometric rate for steep smooth fronts), and linear trial spaces require tens or hundreds of global modes merely to avoid severe Gibbs oscillations~\cite{lanthaler2022nonlinear}.

In recent years, non-intrusive Scientific Machine Learning (SciML) architectures have sought to overcome both intrusiveness and the Kolmogorov barrier by learning either the solution operator or a low-dimensional latent dynamical system directly from FOM trajectories. Broadly, current neural architectures for spatio-temporal PDEs fall into two paradigms:
\begin{enumerate}[leftmargin=2em]
  \item \textbf{Separable modal and operator architectures}, which expand the output field on a fixed or learned spatial basis $\hat{y}(\mathbf{x},t) = \bar{y}(\mathbf{x}) + \sum_{k=1}^K c_k(t)\,\varphi_k(\mathbf{x})$. This family includes PCA-Net~\cite{bhattacharya2021model}, Spectral Neural Operators (SNO), Fourier Neural Operators (FNO)~\cite{li2021fourier} (which alternate linear Fourier truncation with pointwise nonlinear activations on uniform grids), and DeepONets~\cite{lu2021learning} (where the branch network outputs coefficients $c_k$ and the trunk network provides learned basis functions $\tau_k(\mathbf{x})$). As proven by Lanthaler et al.~\cite{lanthaler2022nonlinear}, any operator architecture whose final reconstruction step is a linear combination of $K$ fixed spatial functions remains subject to the linear $K$-width lower bound, regardless of how deep or expressive the coefficient map $c_k$ is.
  \item \textbf{Nonlinear manifold and implicit coordinate architectures}, which parameterise the solution manifold as the image of a low-dimensional latent vector $\mathbf{s}(t)\in\R^{d_s}$ through a nonlinear decoder. Convolutional autoencoder ROMs (DL-ROMs~\cite{lee2020model,fresca2021comprehensive}) compress grid snapshots through an encoder $\mathcal{E}$, evolve $\mathbf{s}(t)$ in latent space, and map back to the grid via a deep convolutional decoder $\mathcal{D}(\mathbf{s}(t))$. More recently, Regazzoni, Pagani, Salvador, Ded\`e, and Quarteroni~\cite{regazzoni2024learning} introduced \emph{Latent Dynamics Networks} (LDNets), which dispense with the encoder altogether and replace the grid-bound decoder with a meshless, implicit neural representation $\NNrec(\mathbf{z}(t),\mathbf{x})$ that takes the spatial query coordinate $\mathbf{x}\in\Omega$ directly as an input alongside the latent state $\mathbf{s}(t)$.
\end{enumerate}

\subsection{The Architecture of Latent Dynamics Networks}

An LDNet models a time-dependent physical system driven by an external time-varying input signal $\mathbf{u}(t)\in\R^{d_u}$ through two coupled multilayer perceptrons (MLPs) trained jointly end to end:
\begin{equation}
  \begin{dcases}
    \dot{\mathbf{s}}(t) = \NNdyn\!\big(\mathbf{s}(t), \mathbf{u}(t)\big), & t \in (0, T], \qquad \mathbf{s}(0) = \mathbf{0}, \\[0.4em]
    \hat{y}(\mathbf{x},t) = \NNrec\!\big(\mathbf{z}(t), \mathbf{x}\big), & \mathbf{x} \in \Omega, \quad t \in [0, T],
  \end{dcases}
  \label{eq:ldnet_continuous}
\end{equation}
where $\mathbf{z}(t) = \mathbf{s}(t)$ (or $\mathbf{z}(t) = [\mathbf{s}(t); \mathbf{u}(t)]$ when instantaneous boundary or forcing inputs act directly on the spatial field). Because $\mathbf{s}(0)=\mathbf{0}$ is fixed at the rest state and no encoder is required to project high-dimensional spatial fields onto the latent space, the entire trajectory $\mathbf{s}(t)$ is determined causally by integrating $\NNdyn$ forward in time from the known input signal $\mathbf{u}(t)$. Furthermore, because $\mathbf{x}$ enters $\NNrec$ as a continuous coordinate, LDNets are completely meshless: during training, the loss function can be evaluated on a small random subset of $N_x \ll N_h$ spatial points per snapshot, drastically reducing memory footprint while still allowing evaluation on the full computational mesh at test time.

However, the coupling in~\eqref{eq:ldnet_continuous} also introduces a fundamental architectural and computational asymmetry between the two sub-networks:
\begin{itemize}[leftmargin=2em]
  \item The \textbf{dynamics network} $\NNdyn$ operates purely in the low-dimensional latent space $\R^{d_s + d_u} \to \R^{d_s}$. At every time step $t_i$ ($i=0,\dots,N_t-1$), it is evaluated \emph{once per trajectory}, sequentially in time.
  \item The \textbf{reconstructor network} $\NNrec$ maps $\R^{\dim(\mathbf{z}) + d_x} \to \R^{n_{\mathrm{comp}}}$. At every time step $t_i$, it must be evaluated \emph{independently at every spatial query point} $\mathbf{x}_j$ ($j=1,\dots,N_x$).
\end{itemize}
In the original LDNet paper~\cite{regazzoni2024learning}, $\NNrec$ contains $85\%$ to $96\%$ of the total trainable parameters (e.g., $1970$ out of $2055$ parameters in the 2D Navier--Stokes cavity flow, and $1480$ out of $1708$ parameters in the 1D cardiac electrophysiology problem), and, because it is evaluated $N_x$ times per time step, it accounts for over $99\%$ of the floating-point operations (FLOPs) in every forward and backward pass!

\subsection{The Two Research Questions: Reconstructor Nonlinearity and Computational Scalability}

The dominance of $\NNrec$ in both parameter count and arithmetic cost motivates the two complementary investigations reported in this paper, denoted \texttt{RQ2} and \texttt{RQ2-bis}:

\paragraph{Part I (\texttt{RQ2}): How much nonlinearity does the LDNet reconstructor actually need?}
In~\eqref{eq:ldnet_continuous}, $\NNrec(\mathbf{z}(t),\mathbf{x})$ performs two distinct tasks simultaneously: it defines the spatial basis functions over $\Omega$, and it maps the latent coordinates $\mathbf{z}(t)$ to the amplitudes of those spatial structures. If we replace $\NNrec(\mathbf{z},\mathbf{x})$ by a modal expansion on a fixed orthonormal basis $\{\varphi_k(\mathbf{x})\}_{k=1}^K$ of $L^2(\Omega)$, we can decompose reconstruction into a hierarchy of four variants:
\begin{enumerate}[label=\textbf{L\arabic*},start=0,leftmargin=2.5em]
  \item \textbf{(Original Pointwise MLP):} $\hat{y}(\mathbf{x},t) = \sigma_{\mathrm{out}}\!\big(\NNrec(\mathbf{z}(t),\mathbf{x})\big)$, nonlinear and non-separable in $(\mathbf{z},\mathbf{x})$.
  \item \textbf{(Parameter-Free Modal Identity):} $\hat{y}(\mathbf{x},t) = \bar{y}(\mathbf{x}) + \sum_{k=1}^{d_s} s_k(t)\,\varphi_k(\mathbf{x})$, identifying each latent state $s_k$ directly with the $k$-th modal amplitude ($K = d_s$, zero decoder parameters).
  \item \textbf{(Affine Modal Synthesis):} $\hat{y}(\mathbf{x},t) = \bar{y}(\mathbf{x}) + \sum_{k=1}^K \big(W\mathbf{z}(t) + \mathbf{b}\big)_k \varphi_k(\mathbf{x})$, decoupling the basis size $K$ from the latent dimension $d_s$ while remaining strictly linear in $\mathbf{z}$ and in space.
  \item \textbf{(Nonlinear Coefficient Map on Fixed Basis):} $\hat{y}(\mathbf{x},t) = \bar{y}(\mathbf{x}) + \sum_{k=1}^K g_\theta(\mathbf{z}(t))_k \varphi_k(\mathbf{x})$, where $g_\theta$ is an MLP parameter-matched to $\mathrm{L0}$ (analogous to a DeepONet with a fixed orthonormal trunk).
\end{enumerate}
By evaluating $\mathrm{L0}$--$\mathrm{L3}$ across two training regimes---\textbf{R1 (frozen latents)}, where $\mathbf{s}(t)$ is frozen from a trained $\mathrm{L0}$ model and only the new reconstructor is fitted, and \textbf{R2 (end-to-end)}, where $\NNdyn$ and the reconstructor are trained jointly from scratch---we isolate whether the nonlinearity of $\NNrec$ is required to break the Kolmogorov $n$-width barrier in space, to untangle curved latent coordinates in state space, or merely to absorb coordinate reparametrisations that an affine modal decoder ($\mathrm{L2}$) could eliminate altogether.

\paragraph{Part II (\texttt{RQ2-bis}): What is the computational trade-off between error, network capacity, BFGS iterations, and FLOPs?}
The \texttt{RQ2} ablation across 816 runs held the dynamics network $\NNdyn$ fixed at its canonical architecture and used a reduced optimization budget (200 Adam + 500 BFGS epochs in Cases~B and~C) to remain computationally tractable on a single workstation. While that design cleanly separated the reconstructor families, it raised three crucial questions for numerical PDE practitioners:
\begin{enumerate}[leftmargin=2em]
  \item When an affine modal reconstructor ($\mathrm{L2}$) has a tiny projection floor $\varepsilon_{\mathrm{proj}}(K)$ (as in the 2D Navier--Stokes cavity flow, where 17 POD modes already reach $\NRMSE \approx 3.5\cdot 10^{-4}$), is the remaining end-to-end error caused by the limited capacity of the small canonical dynamics network $\NNdyn$?
  \item How much of the observed error in $\mathrm{L0}$ and $\mathrm{L2}$ was an artifact of stopping BFGS at 500 or 1,000 iterations rather than running to $10{,}000$ quasi-Newton iterations?
  \item Once $\NNrec$ is replaced by $\mathrm{L2}$, how do the test error, training FLOPs, multi-core wall-clock time, and physical diagnostics (incompressibility, vorticity, front arrival time) scale as we sweep dynamics width $w$, dynamics depth $L$, latent dimension $n$, basis size $K$, and BFGS iteration budget on a high-performance computing (HPC) cluster?
\end{enumerate}
To answer these questions rigorously, \texttt{RQ2-bis} drops the Adam warm-up phase (0 Adam epochs) so that a single quasi-Newton optimization trajectory with passive checkpoints at $k \in \{25, 50, 100, 200, 500, 1000, 2000, 5000, 10000\}$ completed BFGS iterations traces out the entire surface $\mathrm{error}(P, \text{iterations}, \FLOPs)$. Following a 126-run laptop pilot at $1{,}000$ iterations that revealed all convergent $\mathrm{L0}$ and wide $\mathrm{L2}$ models were still dropping steeply at iteration $1{,}000$, we executed the full $10{,}000$-iteration campaign (126 runs across Cases~B and~C) plus dedicated exclusive-node thread-scaling and evaluation-timing benchmarks on the \textbf{CINECA Leonardo supercomputer} (\texttt{dcgp\_usr\_prod} partition, dual-socket Intel Xeon Platinum 8480+ nodes with 112 physical cores per node).

\subsection{Summary of Main Contributions and Paper Organization}

The contributions of this work to the NMPDE and SciML literature are fivefold:
\begin{enumerate}[leftmargin=2em]
  \item \textbf{Theoretical rank bounds for coupled latent-dynamics/reconstructor architectures (Section~\ref{sec:math}):} In addition to formalising the \emph{reconstructor rank bound} $\operatorname{rank}(\hat{y}-\bar{y}) \le \min(\dim\mathbf{z}, K)$ for affine modal decoders ($\mathrm{L2}$), we prove a new \emph{dynamics hidden-width bottleneck bound} (Proposition~\ref{prop:dyn_rank_bound}): when $\NNdyn$ has last hidden layer width $w_{\mathrm{dyn}} < n$, Forward Euler integration from $\mathbf{s}(0)=\mathbf{0}$ confines both the trajectory matrix $\mathbf{s}(t)$ and the state-transition Jacobian $\partial\dot{\mathbf{s}}/\partial\mathbf{s}$ to a subspace of dimension at most $\min(n, w_{\mathrm{dyn}}+1)$ (and strictly $\min(n, w_{\mathrm{dyn}})$ under exact rest-state equilibrium subtraction). This proves that increasing the latent dimension $n$ beyond the hidden width of $\NNdyn$ is structurally vacuous unless $\NNdyn$ is widened simultaneously.
  \item \textbf{Complete reconstructor nonlinearity taxonomy across PDE regimes (Section~\ref{sec:rq2}):} Across 816 runs in \texttt{RQ2}, we demonstrate that whether $\NNrec$ can be replaced by a linear reconstructor is governed by the Kolmogorov $n$-width decay of the solution manifold rather than the nonlinearity of the PDE itself. On smooth periodic advection--diffusion--reaction (Case~A) and 2D Navier--Stokes cavity flow (Case~B), linear and fixed-basis decoders ($\mathrm{L2}$, $\mathrm{L3}$) trained end to end match or outperform $\mathrm{L0}$, whereas frozen-latent ($\mathrm{R1}$) experiments prove that $\mathrm{L0}$ learns curved, non-modal latent coordinates even on strictly linear PDEs. Conversely, on 1D cardiac electrophysiology with sharp travelling fronts (Case~C), only the non-separable pointwise MLP $\mathrm{L0}$ breaks the Kolmogorov barrier, collapsing the error by an order of magnitude precisely when the latent dimension crosses the intrinsic manifold threshold $d_s = 3$.
  \item \textbf{Asymptotic capacity laws and structural rank collapse at 10,000 BFGS iterations (Section~\ref{sec:rq2bis}):} Using the CINECA Leonardo $10{,}000$-iteration sweeps, we show that power-law capacity scaling $E(P_{\mathrm{dyn}}) = E_\infty + a\,P_{\mathrm{dyn}}^{-\alpha}$ emerges if and only if neither the reconstructor nor the latent dimension is the active bottleneck. In Case~B, $\mathrm{L0}$ at $n^*=1$ is completely flat across a $62\times$ range of dynamics parameters ($\alpha = 0.12$), whereas replacing $\mathrm{L0}$ with the linear POD reconstructor $\mathrm{L2}$ at $n=10, K=65$ unlocks a clean power law ($\alpha = 1.11$, 95\% CI $[1.04, 1.23]$) that drives test NRMSE down to $2.32\cdot 10^{-3}$---$2.5\times$ more accurate than the best $\mathrm{L0}$ configuration ($5.86\cdot 10^{-3}$), with $4.7\times$ lower vorticity error, $9.5\times$ lower velocity direction error, and exact preservation of the FOM's discrete divergence floor. In Case~C, $\mathrm{L0}$ exhibits steep capacity scaling ($\alpha = 3.22$) where increasing dynamics depth ($L=3$, NRMSE $7.83\cdot 10^{-3}$, front arrival-time MAE $0.54$~ms) beats quadrupling dynamics width ($w=32$, NRMSE $8.57\cdot 10^{-3}$) at less than half the parameter count, while wide $\mathrm{L2}$ models overfit high-degree Legendre oscillations between the 20 subsampled spatial training points after iteration $1{,}000$.
  \item \textbf{Analytical and hardware-measured FLOPs/error scalability (Section~\ref{sec:flops}):} We derive the exact forward/backward operation count $\FLOPs_{\mathrm{eval}}$ and confront it with exclusive-node Pass-T wall-clock timings across $p \in \{1,2,4,8,16,32,64,112\}$ CPU threads on CINECA Leonardo. We show that in Case~B, $\mathrm{L2}$ reaches the best $\mathrm{L0}$ accuracy using $56.3\times$ fewer FLOPs and $10.3\times$ less exclusive-node wall-time ($347$~s vs.\ $3556$~s), and reaches $1.5\times$ the best $\mathrm{L0}$ accuracy using $167\times$ fewer FLOPs and $11.9\times$ less wall-time. At the same time, we explain the $3.0\times$--$3.1\times$ discrepancy between FLOP ratios and wall-time ratios via hardware arithmetic intensity: $\mathrm{L0}$ executes dense batched matrix multiplications at $\sim 60.9$~GFLOP/s (Case~B) and $\sim 32.4$~GFLOP/s (Case~C) on 16 Xeon cores, whereas $\mathrm{L2}$ strips away the spatial batch work and exposes the sequential $T$-step Forward Euler loop, dropping sustained throughput to $4.4$--$20.2$~GFLOP/s and capping multi-core parallel speedup via Amdahl's law ($23.5\times$ at 112 threads for $N_t=101$ in Case~B vs.\ $8.6\times$ for $N_t=501$ in Case~C).
  \item \textbf{Methodological and numerical guidelines for SciML (Section~\ref{sec:discussion}):} We synthesise nine concrete design rules for latent dynamical operator learning, covering the coupling restriction $w_{\mathrm{dyn}} \ge n$, the necessity of multi-thousand-iteration quasi-Newton budgets, the initial-state handicap of parameter-free centred modal expansions ($\mathrm{L1}$), spatial subsampling aliasing in high-degree polynomial bases, and cross-node AVX-512 floating-point non-determinism in long BFGS trajectories.
\end{enumerate}

The remainder of this paper is organised as follows. Section~\ref{sec:math} presents the mathematical formulation of LDNets, the $\mathrm{L0}$--$\mathrm{L3}$ reconstructor hierarchy, the discrete inner product and orthonormalisation procedures, and the proofs of the reconstructor and dynamics bottleneck rank bounds. Section~\ref{sec:problems} describes the three benchmark PDE problems and their a priori singular-value spectra and projection floors before any training. Section~\ref{sec:rq2} reports the complete results of the \texttt{RQ2} reconstructor nonlinearity ablation across Cases~A, B, and~C. Section~\ref{sec:rq2bis} presents the \texttt{RQ2-bis} $10{,}000$-iteration capacity, depth, latent-dimension, and iteration-budget scalability study on the CINECA Leonardo supercomputer. Section~\ref{sec:flops} provides a dedicated, self-contained analysis of computational scalability in terms of FLOPs, hardware throughput, multi-core strong scaling, and error. Section~\ref{sec:discussion} consolidates the hypotheses scorecards, practical engineering rules, limitations, and future directions, while Appendix~\ref{sec:appendix} provides supplementary verification tables and reproducibility protocols.
